\section{Final Questions, Answered Directly}
\label{sec:final-questions}

\begin{enumerate}[itemsep=0.4em]

\item \textbf{What should physically be identical between E2B and Gateway hardware?}
MCU and clock tree; the LAN8650/LAN8651 network block and its coupling/termination; power input protection, regulation, and voltage/current monitoring; reset and watchdog circuitry; the SWD debug/programming header and pinout; the timestamping resource; the manufacturing/test procedure; and the PCB stackup and design rules (Section~\ref{sec:common-platform}).

\item \textbf{What should differ only through component population or configuration?}
Which of the two reserved footprint areas is stuffed (peripheral-interface section vs. legacy-bus section), and which of two thin firmware role modules is compiled into the shared base image. Nothing else should differ between an E2B build and a Gateway build of Rev~B.

\item \textbf{What functions require separate interface circuitry?}
Peripheral conditioning (GPIO/PWM/ADC/UART/SPI signal conditioning specific to sensors/actuators) and legacy-bus transceivers (CAN-FD, optional LIN) plus their hardware TX-enable gating -- these are electrically distinct front ends that cannot share silicon, only a common host MCU and common power/debug infrastructure.

\item \textbf{Baseboard + daughtercard, or closely related PCB variants?}
Neither, for Rev~A/B -- one common PCB with selective population (Section~\ref{sec:common-platform}), because it is the lowest-risk option for September--November iteration speed. Closely related PCB variants (not a daughtercard connector) for Rev~C, once Rev~B has resolved footprint/pin conflicts. A true baseboard+daughtercard split is explicitly rejected for this program: it adds a board-to-board connector and a second layout risk surface neither the schedule nor the requirement (``prove the common core, not build a modular product line'') justifies.

\item \textbf{What exact Microchip development hardware should Alfred purchase immediately?}
Three MikroE Two-Wire ETH Click (MIKROE-5543, LAN8651B1), one MikroE Two-Wire ETH~3 Click (MIKROE-6550) as a revision cross-check, one Microchip EVB-LAN8670-USB as an independent reference node, and loose LAN8650/LAN8651 samples plus an STM32G4 Nucleo board for Rev~A/B schematic and firmware work (Section~\ref{sec:lan866x}). LAN8660/LAN8661 dedicated evaluation hardware is not purchased on this timeline -- no Microchip EVB for LAN8660 was found in current documentation, and the part is a secondary, time-boxed investigation, not a critical-path item.

\item \textbf{What does Rev A teach us?}
Whether the OA-TC6 SPI link comes up cleanly on Alfred's chosen MCU using the open reference driver; measured PLCA transmit-opportunity timing with 2--3 real nodes; real T1S supply current at idle and under sustained TX; link sensitivity to the lab's actual bench cable/termination setup; and, from its secondary section, a first-hand, evidence-based answer on whether the LAN8660 MCU-less endpoint is worth further product investment at all (Section~\ref{sec:rev-a}).

\item \textbf{What architecture must Rev B prove?}
That one populated-selectively PCB, one MCU, one network block, one power architecture, and one firmware base image can be built up as an E2B node and, separately, as a Gateway node, each independently meeting its own minimum demonstration (Sections~\ref{sec:arch-a-demo}, \ref{sec:arch-b}) -- not that either role works in isolation on bespoke hardware (Section~\ref{sec:rev-b}).

\item \textbf{What makes Rev C credible enough to show an automotive customer?}
Locking connectors, demo-appropriate mounting and cabling, self-explanatory labeling, a documented and repeatable configuration procedure, and BOM/testability quality consistent with a serious engineering program -- explicitly not AEC qualification, production EMC, ISO~26262 hardware, or environmental qualification, none of which the exit gate requires (Section~\ref{sec:rev-c}).

\item \textbf{What minimum experiment validates the clean-sheet architecture?}
Two Alfred nodes in E2B mode on a multidrop 10BASE-T1S segment with a Clicker~4, carrying one real sensor input and one real actuator command, with PLCA behavior, latency, jitter, disconnect/reconnect, and power-cycle recovery all measured, not asserted (Section~\ref{sec:arch-a-demo}).

\item \textbf{What minimum experiment validates the brownfield architecture?}
An Alfred Gateway on a controlled bench CAN network, first in Mode~1 (proving passive capture/timestamp/mirror onto the T1S pseudo-bus with verified zero TX), then in Mode~2 (a Clicker-issued command reaching a physical device through the Gateway and CAN, with the device's response observed and mirrored back, and round-trip latency measured) (Section~\ref{sec:arch-b}).

\item \textbf{Can all of this realistically be completed by November 30?}
The core thesis (Sections~\ref{sec:exit-gate}'s two checklists, run on Rev~B) is realistic on this schedule \emph{if} Rev~A and Rev~B each ship and arrive on their first fabrication pass and if LIN and Rev~C are treated as genuinely conditional rather than assumed. Three clean PCB spins plus full customer-demo polish (Rev~C, enclosure, cabling, both buses including LIN) in twelve weeks, with no fabrication or bring-up surprises, is \emph{not} a credible planning assumption for a small team -- it is an optimistic path, not the base case.

\item \textbf{If not, what must be removed while preserving the central Alfred thesis?}
Follow Section~\ref{sec:feature-cuts} exactly: drop LIN, drop Rev~C, drop any protocol beyond CAN/CAN-FD, drop enclosure/cosmetic polish, and drop the LAN8660 secondary investigation before touching any item on the ``protect at all costs'' list. The central thesis survives fully intact on characterized Rev~B hardware alone: one common node, populated two ways, each independently meeting its minimum demonstration with quantitative data. That is the irreducible core; everything else in this plan is schedule margin spent making the proof more complete or more presentable, not the proof itself.

\end{enumerate}
